\documentclass[10pt,a4paper]{amsart}
\usepackage[margin=19mm]{geometry}
\usepackage{amsmath,amssymb,mathtools}
\usepackage[scr=rsfs,cal=euler]{mathalfa}
\usepackage{xfrac}
\usepackage[unicode,hidelinks,bookmarks=false]{hyperref}
\newcommand{\ie}{i.e.\ }
\newcommand{\cf}{cf.\ }
\newcommand{\resp}{respectively\ }
\newcommand{\Subsection}{\subsection}
\providecommand{\nobreakdash}{\nobreak\hbox{-}}
\setlength{\emergencystretch}{2em}
\multlinegap0pt
\usepackage{amssymb}
\newtheorem{thm}{Theorem}
\newtheorem{lem}{Lemma}
\newcommand{\diam}{\operatorname{Diam}}
\newcommand{\vol}{\operatorname{vol}}
\title{A note on the distribution of Neumann eigenvalues of the Laplacian on a Euclidean convex domain}
\author{Kei Funano}
\date{Source: arXiv:2603.16135v1 (CC BY 4.0)}
\begin{document}
\maketitle

\noindent\emph{Source and licence.} A note on the distribution of Neumann eigenvalues of the Laplacian on a Euclidean convex domain. Version arXiv:2603.16135v1 (CC BY 4.0), distributed by arXiv under the Creative Commons Attribution 4.0 International licence, http://creativecommons.org/licenses/by/4.0/, as recorded in the arXiv licence metadata for this exact version and shown on the abstract page https://arxiv.org/abs/2603.16135v1. Full licence notice: \texttt{licenses/arxiv-2603.16135-CC-BY-4.0.txt}.\par\smallskip

\noindent\textit{Editorial scope.} Counted unit: the article's lem Mlem1 (source0.tex lines 200--204). The article's complete original proof of it is delivered with it (source0.tex lines 229--277, one \texttt{proof} environment of 49 lines, which the article places away from the statement). No mathematics is written, repaired or re-derived: the statement and the proof are the article's own lines, in the article's own notation, and no retained line is substituted or re-flowed.  Retained uncounted as the source-local context the proof needs: lines 176--180.  Nothing was omitted from any retained span, so the package carries no omission record.  The retained lines cite 1 works, whose entries are transcribed verbatim from the article's own bibliography source in the e-print and delivered with short editorial labels recorded in the item's reference mapping.  No retained line contains a figure environment, an image-inclusion command or drawing code, so no image is delivered and the package declares no asset dependency.  Editorial formatting only, in addition: an AMS \texttt{amsart} preamble that restates the article's own declarations and macros used by the retained lines, namely the environments \texttt{thm}, \texttt{lem} on separate counters, together with the standard packages those lines need.  The retained environments print in this document as Theorem 1, Lemma 1 in order of appearance, whereas the article prints the same items with the numbering of its own full document.  The complete native source of the article as distributed in the arXiv e-print for this exact version is bundled unchanged as \texttt{sources/196519/source0.tex}.
\begin{thm}[{Kr\"{o}ger, \cite[Corollary 2]{Kr2}}]\label{Kr ineq2}Let $\Omega$ be a (not necessarily convex) bounded domain in $\mathbb{R}^n$ with Lipschitz boundary. Then for any natural number $k$ we have
\begin{align*}
    \mu_k(\Omega)\leq (2\pi)^2\Big(\frac{n+2}{2}\Big)^{\frac{2}{n}}\Big(\frac{k}{\omega_n \vol \Omega}\Big)^{\frac{2}{n}}\lesssim \Big(\frac{k}{\omega_n \vol \Omega}\Big)^{\frac{2}{n}},
\end{align*}where $\omega_n$ is the volume of a unit ball in $\mathbb{R}^n$.
\end{thm}

\begin{lem}\label{Mlem1}There exists a universal numeric constant $c>0$ satisfying the following. For any orthotope $\Omega$ in $\mathbb{R}^n$ and for any $k\geq l$ there exists a convex partition $\{\Omega_i\}_{i=1}^{l'}$ of $\Omega$ such that $l'\leq l$ and 
    \begin{align*}
        \diam \Omega_i \leq cn^{\frac{3}{2}}\frac{k}{l\sqrt{\mu_k(\Omega)}}.
    \end{align*}
\end{lem}

\begin{proof}[Proof of Lemma \ref{Mlem1}]
    Take a universal constant $C>0$ such that 
    \begin{align*}
         C\geq \Big((2\pi)^n\cdot \frac{n+2}{2} \cdot \frac{1}{(\omega_n)^2}\Big)^{\frac{1}{n}}\cdot \frac{1}{n}\sim \text{ const. }
    \end{align*}for any $n\geq 2$. Also take a universal constant $c>0$ such that
    \begin{align*}
        c\geq 4C \text{ and }c>\frac{2\sqrt{n}}{\sqrt{n}-1} \text{ for any }n\geq 2.
    \end{align*}

    
    For such $c$ we prove the lemma by induction on $n$. The case where $n=1$ is trivial. 

    Suppose that the claim of the lemma holds for any orthotope $\Omega$ in $\mathbb{R}^{n-1}$, that is, suppose that for such $\Omega$ we can always find a convex partition $\{\Omega_i\}_{i=1}^{l'}$ of $\Omega$ such that $l'\leq l$ and $\diam \Omega_i\leq c(n-1)^{\frac{3}{2}}\frac{k}{l\sqrt{\mu_k(\Omega)}}$.

    Let $\Omega$ be an orthotope in $\mathbb{R}^n$. After choosing the coordinate axes appropriately we may assume that $\Omega=[-a_1,a_1]\times [-a_2,a_2]\times \cdots \times [-a_n,a_n]$ and $a_1\leq a_2\leq \cdots \leq a_n$.


    
    Put $r:=\frac{Cnk}{l\sqrt{\mu_k(\Omega)}}$. We first consider the case where $a_1\leq 2r$. Set $\Omega':=[-a_2,a_2]\times [-a_3,a_3]\times \cdots \times [-a_n,a_n]$. By the assumption of the induction there exists a convex partition $\{\Omega_i'\}_{i=1}^{l'}$ of $\Omega'$ such that $l'\leq l$ and $\diam \Omega_i'\leq c(n-1)^{\frac{3}{2}}\frac{k}{l\sqrt{\mu_k(\Omega')}}$. Note that $\mu_k(\Omega)\leq \mu_k(\Omega')$ (This follows from the concrete expression of eigenvalues of an orthotope). Thus we get $\diam \Omega_i'\leq c(n-1)^{\frac{3}{2}}\frac{k}{l\sqrt{\mu_k(\Omega)}}$. Setting $\Omega_i:=[-a_1,a_1]\times \Omega_i'$ we see that $\{\Omega_i\}_{i=1}^{l'}$ covers $\Omega$. We also have 
    \begin{align*}
        \diam \Omega_i\leq \sqrt{4^2C^2n^2+c^2(n-1)^3}\frac{k}{l\sqrt{\mu_k(\Omega)}}\leq cn^{\frac{3}{2}}\frac{k}{l\sqrt{\mu_k(\Omega)}},
    \end{align*}where the second inequality follows from our choice of the constants $c$ and $C$. Therefore $\{\Omega_i\}_{i=1}^{l'}$ is a desired convex partition of $\Omega$.
    

    Let us consider the case where $a_1>2r$. In this case we put 
    \begin{align*}
        \widetilde{\Omega}:=\{x\in \Omega \mid d(x,\partial \Omega)\geq r\},
    \end{align*}where 
    \begin{align*}
        d(x,\partial \widetilde{\Omega}):=\inf \{|x-y| \mid y\in \partial \widetilde{\Omega}\}.
    \end{align*}
    Note that $\widetilde{\Omega}$ is a nonempty (noncollapsed) orthotope. Let $\{x_i\}_{i=1}^{l'}$ be a maximal $(2r)$-separated points in $\widetilde{\Omega}$. Since $\{B(x_i,r)\}_{i=1}^{l'}$ are disjoint balls in $\Omega$, we have
    \begin{align*}
       \frac{l'\omega_n(Cnk)^n}{l^n\mu_k(\Omega)^{\frac{n}{2}}}= l'\omega_n r^n =\sum_{i=1}^{l'} \vol B(x_i,r)\leq \vol \Omega.
    \end{align*}
    Applying Theorem \ref{Kr ineq2} to the above inequality we have
    \begin{align*}
        l'\leq (2\pi)^n\frac{n+2}{2}\cdot \frac{1}{(\omega_n)^2}\cdot \frac{l^{n-1}}{C^n n^nk^{n-1}}\cdot l\leq \frac{l^{n-1}}{k^{n-1}}\cdot l \leq l,
    \end{align*}where the second inequality follows from our choice of $C$ and $k\geq l$ implies the last inequality. 
   

    The maximality of the points $\{x_i\}_{i=1}^{l'}$ gives $\widetilde{\Omega}\subseteq \bigcup_{i=1}^{l'}B(x_i,2r)$. Also observe that the $\sqrt{n}r$-neighbourhood of $\widetilde{\Omega}$ covers $\Omega$ (note that the $r$-neighbourhood of $\widetilde{\Omega}$ does not cover $\Omega$). We thereby have $\Omega\subseteq \bigcup_{i=1}^{l'}B(x_i,(2+\sqrt{n})r)$. Therefore if we introduce the \emph{Voronoi cells}
    \begin{align*}
        \Omega_i:=\{x\in \Omega \mid |x-x_i|\leq |x-x_j| \text{ for all }j\neq i\},
    \end{align*}then $\{\Omega_i\}_{i=1}^{l'}$ is a convex partition of $\Omega$. Since $\Omega_i\subseteq B(x_i,(2+\sqrt{n})r)$ our choice of $c$ implies
    \begin{align*}
        \diam \Omega_i\leq 2(2+\sqrt{n})r\leq c n^{\frac{3}{2}}\frac{k}{l\sqrt{\mu_k(\Omega)}}.
    \end{align*}Therefore we obtain the desired convex partition $\{\Omega_i\}_{i=1}^{l'}$ of $\Omega$. This completes the proof.
\end{proof}

\begin{thebibliography}{99}
\bibitem[Kr2]{Kr2} P.~Kr\"{o}ger, \textit{Upper bounds for the Neumann eigenvalues on a bounded domain in Euclidean space}. J. Funct. Anal.106(1992), no.2, 353--357. %
\end{thebibliography}
\end{document}
